% book模板
% book模板.tex

\documentclass[12pt,UTF8]{ctexbook}

% 设置纸张信息。
\input{../../../../../common/page}

% 设置字体，并解决显示难检字问题。
\xeCJKsetup{AutoFallBack=true}
% 注意：ctexbook类已默认设置SimSun为CJKrmdefault
% 以下设置用于确保字体回退和扩展字体可用
% 仅设置扩展字体以避免与默认设置冲突
\setCJKfamilyfont{hei}{SimHei}
\setCJKfamilyfont{kai}{KaiTi}

% 目录 chapter 级别加点（.）。
\usepackage{titletoc}
\titlecontents{chapter}[0pt]{\vspace{3mm}\bf\addvspace{2pt}\filright}{\contentspush{\thecontentslabel\hspace{0.8em}}}{}{\titlerule*[8pt]{.}\contentspage}

% 支持目录点击跳转
\usepackage[colorlinks,linkcolor=blue,citecolor=blue,urlcolor=blue]{hyperref}

% 设置 part 和 chapter 标题格式。
\ctexset{
	part/name= {第,卷},
	part/number={\chinese{part}},
	chapter/name={第,篇},
	chapter/number={\arabic{chapter}}
}

% 图片相关设置。
\usepackage{graphicx}
\graphicspath{{Images/}}

% 设置署名格式。
\newenvironment{shuming}{\hfill\zihao{4}}

% 注脚每页重新编号，避免编号过大。
\usepackage[perpage]{footmisc}

% 设置古文原文格式。
\newenvironment{yuanwen}{\bfseries\zihao{4}}

% 设置段落间距
\setlength{\parskip}{0.8em plus 0.2em minus 0.1em}

\usepackage{float}

\usepackage{caption}
\usepackage{tabularx}
\usepackage{array}
\usepackage{hyperref}
\usepackage{booktabs}
\usepackage{longtable}

% 列表格式支持，列表项向右偏移。
\usepackage{enumitem}

\title{\heiti\zihao{0} 题目}
\author{WangFei}
\date{\today}

\begin{document}

\maketitle
\tableofcontents

\frontmatter
\chapter{前言、序言}

信号发生器是电子测量中的"激励源"。在测量系统中，仪器大体分为两类：一类用于观察和读取被测量，如电压表、示波器、频谱分析仪；另一类用于向被测电路提供已知的、参数可控的信号，使被测电路进入预定的工作状态，从而能够观察其响应。后者即信号发生器，或称信号源。没有可信的激励，就无法定量地考察放大器的增益与频响、谐振回路的调谐与选择性、接收机的灵敏度与通带，因而信号源与测量仪表同为电子实验室的两块基石。

一台通用信号发生器在结构上大致可分为三个环节。第一个环节是\textbf{频率产生}，即用振荡器得到所需频率的信号；早期仪器多采用 LC 振荡器或 RC 振荡器并以可变电容、可变电阻调节频率，现代仪器则普遍采用石英晶体基准加锁相环或直接数字合成的频率合成方式，以获得高得多的频率准确度和分辨力。第二个环节是\textbf{电平控制}，即用放大器、宽带衰减器和自动电平控制电路，使输出幅度在宽范围内连续可调且在整个频段内保持平坦。第三个环节是\textbf{调制}，即按需要对载波施加调幅、调频、调相或脉冲调制，以模拟实际通信信号。三个环节之后再经输出级和输出接口把信号送出，输出级同时承担阻抗变换与隔离，使外接负载的变化不至于牵动振荡器的频率。

按用途与工作频段，信号源大致有以下几类。\textbf{低频（音频）信号发生器}产生数赫兹至数百千赫的正弦信号，波形失真小，用于音频通道与低频放大器的测量；\textbf{函数发生器}可产生正弦波、方波、三角波、锯齿波并可调直流偏置与占空比，通用性强，常用于数字与模拟电路调试；\textbf{任意波形发生器}以存储波形数据后经数模转换输出的方式工作，可复现任意给定波形，适合模拟实际信号与故障波形；\textbf{高频与射频信号发生器}覆盖数百千赫至数吉赫，具备精确的电平控制与调制功能，是接收机测试的主要激励源；此外还有专用于测量幅频特性的\textbf{扫频信号发生器}和专用于时序测量的\textbf{脉冲信号发生器}。本篇所介绍的学生信号发生器同时具备低频与高频两路输出，并可用低频信号对高频信号调幅，属于教学与一般检修用的简易信号源。

评价信号源性能的指标主要有以下几项。\textbf{频率范围}指可输出的频率区间；\textbf{频率准确度}指实际频率与标称（刻度）值之差，常以相对误差表示；\textbf{频率稳定度}指频率随时间、温度和电源电压变化而漂移的程度，短期稳定度影响单次测量，长期稳定度决定校准周期。\textbf{相位噪声}描述载波附近的随机相位起伏，以偏离载波一定频偏处每单位带宽的噪声与载波之比表示，单位为 dBc/Hz，它限制了接收机邻道选择性等指标的可测下限。\textbf{输出电平范围与平坦度}指幅度可调的区间以及在频段内幅度随频率变化的最大偏差。\textbf{谐波与杂散}指输出中除载波外的其他频率成分，谐波为载波频率的整数倍，杂散则包括本振泄漏、电源交流声等非谐波成分，两者都会在测量中造成虚假响应。\textbf{调制指标}包括调幅度范围与调幅失真、频偏范围与调频失真、脉冲调制的开关比与上升时间等。下表汇总了这些指标的物理含义与实用意义。

\begin{table}[H]
	\centering
	\caption{信号发生器主要技术指标及其意义}
	\begin{tabular}{|p{2.6cm}|p{4.6cm}|p{5.4cm}|}
		\hline
		\textbf{指标} & \textbf{含义} & \textbf{对测量的影响} \\ \hline
		频率范围 & 可输出频率的上下限 & 决定可测电路的工作频段 \\ \hline
		频率准确度 & 实际频率与刻度值之差 & 直接影响调谐、对中频等测量的正确性 \\ \hline
		频率稳定度 & 频率随时间与温度的漂移 & 影响长时间测量与窄带电路调整 \\ \hline
		相位噪声 & 载波附近随机相位起伏 & 限制邻道选择性、抗阻塞等测量下限 \\ \hline
		输出电平范围 & 幅度可调的区间 & 决定能否测到接收机灵敏度等小信号指标 \\ \hline
		电平平坦度 & 幅度随频率的变化量 & 影响幅频特性与增益测量的准确度 \\ \hline
		谐波与杂散 & 载波外的其他频率成分 & 造成虚假响应，影响失真与选择性测量 \\ \hline
		调制指标 & 调幅度、频偏及其失真 & 决定解调、鉴频等测量的可信度 \\ \hline
	\end{tabular}
\end{table}

使用信号源时应牢记三条基本原则。其一，注意输出阻抗与负载的关系：多数信号源按额定负载标定输出幅度，负载偏离额定值时实际幅度将改变，射频测量还会因失配产生驻波与幅度起伏。其二，幅度调节应由小到大，尤其在向接收机或高增益放大器注入信号时，过大的输入会使被测电路饱和甚至损坏，得到的现象也无意义。其三，输出端不得短路，也不应把外部直流电压加到输出端，否则容易损坏输出级。现代信号源大多带有程控接口，可在自动测试系统中由程序设定频率、电平与调制参数，从而实现频率步进扫描与批量测量；程控使用时应注意每次参数改变后留出足够的稳定时间。此外，信号源本身也是需要校准的计量器具，其频率、电平与调制指标应定期与更高等级的标准比对；日常使用中应保持通风干燥、避免碰撞，长期存放的仪器宜定期通电运行，以免元件受潮变质。

\mainmatter

\chapter{J2465-2型学生信号发生器 杭州九二无线电厂}
	
\section{概述}

本仪器是在 J2465 型学生信号发生器的基础上改进设计而成的新产品。

仪器的主要改进为：

第一，高频振荡采用新型场效应管电容三点式电路，改善了原仪器存在波形失真大、频率稳定性差、输出电压幅度小等缺点。

第二，高频信号调幅采用场效应管漏极调幅，解决了原仪器调幅时载频发生偏移的问题。

第三，低频振荡器采用 RC 桥式振荡电路，提高了频率稳定度及准确度，并使波形失真度达到小于 1\%。

第四，采用了新颖结构，仪器造型美观大方，使用维修方便。

本仪器主要作为中学物理学生分组实验使用，也可供生产厂、维修单位等作为测试信号源用。

从功能上看，本仪器实际上把两台简易信号源装在一个机壳内，并在两者之间建立了调制联系。低频通道输出五个点频的音频正弦信号，可用于检查音频放大器与扬声器通道、观察放大器的波形与增益，也可作为示波器与音频电压表的实验信号；高频通道输出中波广播频段的正弦信号，可用于检查收音机的输入回路、高频放大与中频放大电路。当把两者接通时，低频信号即成为高频信号的调制信号，仪器输出一个已调幅的高频信号，从而可以完整地检查从天线输入到扬声器输出的整条接收通道。这种"高频加低频、等幅兼调幅"的配置，正是中波超外差收音机检修与教学实验所需的最小配置。

高频通道中 465~kHz 这一频点具有特殊意义。中波超外差收音机的中频普遍取 465~kHz，检修时需要一个准确的中频信号来逐级对准中频变压器，使各级谐振曲线的峰值都落在中频上，从而获得最大增益与合适的通带。本仪器把该频点标注在刻度盘上并单独给出较小的误差指标，就是为此用途。至于等幅波与调幅波的选择：检查高频与中频电路的通路、调整谐振回路时使用等幅波即可，用高频电压表或示波器在各级观察幅度变化；而检查检波器、音频通道以及整机听感时必须使用调幅波，因为只有已调信号才能在检波后还原出音频。

需要说明的是，本仪器的定位是教学与一般检修，其频率与幅度的准确度指标较为宽松。若需要准确知道输出频率，应另接频率计测量；若需要准确知道输出电平，应另接高频电压表或用已知衰减器进行比较。把简易信号源的刻度值直接当作精确量值使用，是初学者常犯的错误。

\section{技术特性}

\begin{enumerate}
	\item \textbf{高频信号}
	\begin{itemize}
		\item 频率范围：450kHz -- 1700kHz 连续可调
		\item 频率刻度误差：\(\le \pm 5\%\)
		\item 中频 465kHz 误差：\(\le \pm 2\%\)
		\item 输出幅度：\(\ge 300\)mV （负载为 300\(\Omega\)）
		\item 调幅度：\(\ge 30\%\)
	\end{itemize}
	\item \textbf{低频信号}
	\begin{itemize}
		\item 频率范围：五个点频 500Hz、1000Hz、1500Hz、2000Hz、2500Hz
		\item 频率误差：\(\le 5\%\)
		\item 输出幅度：\(\ge 800\)mV （负载为 300\(\Omega\)）
		\item 波形失真度：\(< 1\%\)
	\end{itemize}
	\item \textbf{使用条件}
	\begin{itemize}
		\item 温度：\(0^\circ\)C ～ \(+40^\circ\)C
		\item 相对湿度：20\% ～ 90\%
		\item 使用电源：220V \(\pm 10\%\)， 50Hz
		\item 工作时间：连续 8 小时
		\item 外形尺寸：120mm \(\times\) 166mm \(\times\) 70mm
		\item 重量：约 1.5kg
	\end{itemize}
	\item \textbf{附件：}
	\begin{itemize}
		\item 输出电缆：2 根
		\item 使用说明书：1 本
		\item 备用保险丝：0.2A， 2 只
	\end{itemize}
\end{enumerate}

上列各项指标的含义需要逐条理解，才能正确判断仪器能否满足某项测量的要求。\textbf{频率刻度误差}是指按刻度盘读得的频率与实际频率之差相对于实际频率的比值，即
\[
\gamma_f = \frac{f_d - f_x}{f_x} \times 100\%,
\]
式中 $f_d$ 为刻度读数，$f_x$ 为实际频率。该误差包含了刻度印制、指针装配、调谐机构间隙以及回路元件容差等多种因素。$\pm 5\%$ 的量级意味着在 1000~kHz 附近读数可能有数十千赫的偏差，因此本仪器的刻度只能用于粗定频率，用来判断"大致在哪一档"，而不能用于测定被测回路的谐振频率。中频点之所以能给出较小的误差，是因为该点在出厂时经过单独校准。

\textbf{输出幅度}的指标都附有负载条件（$300\,\Omega$），这一点极易被忽略。信号源的输出级具有一定的等效内阻 $R_o$，当外接负载 $R_L$ 时，负载上得到的电压为
\[
U_L = U_o\,\frac{R_L}{R_o + R_L},
\]
式中 $U_o$ 为空载输出电压。负载电阻越小，实际得到的幅度越低；若负载远小于额定值，不仅幅度显著下降，还可能因输出级过载而使波形失真，甚至牵动振荡器频率。因此后文的使用方法中特别要求被测电路的输入电阻不小于 $300\,\Omega$，其道理正在于此。

\textbf{调幅度}描述高频载波被低频信号调制的深浅程度。设已调波包络的最大值为 $U_{\max}$、最小值为 $U_{\min}$，则调幅度定义为
\[
m = \frac{U_{\max} - U_{\min}}{U_{\max} + U_{\min}} \times 100\%.
\]
$m$ 越大，检波后得到的音频信号越强；但当 $m$ 接近或超过 $100\%$ 时包络将出现底部削平的过调制现象，检波输出严重失真。用示波器观察输出包络即可按上式估算调幅度，这也是本仪器可用于教学演示调幅概念的原因。

\textbf{波形失真度}即总谐波失真，定义为各次谐波有效值的方和根与基波有效值之比
\[
THD = \frac{\sqrt{U_2^2 + U_3^2 + U_4^2 + \cdots}}{U_1} \times 100\%.
\]
低频振荡器的失真度小于 $1\%$，说明其波形已相当接近正弦，可以用来测量音频放大器的失真——测量时被测放大器的失真必须明显大于信号源自身的失真，否则读数中将有相当部分来自信号源。

\textbf{使用条件}诸项同样是指标的组成部分而非附带说明。温度与湿度范围之外使用时，各项指标不再保证，尤其是频率稳定度与失真度会明显变坏；电源电压偏离 $220\,\mathrm{V} \pm 10\%$ 时，稳压电路可能失去稳压能力，导致振荡器工作点偏移、幅度变化甚至停振。连续工作时间与外形尺寸、重量则关系到仪器的散热设计与使用场合，成组用于学生实验时应注意课桌上仪器之间留出通风间隔。

\section{控制器作用}

仪器面板上的控制件按功能分为三组：公共部分（电源开关、保险丝、接地端）、低频部分（频率转换开关、幅度调节、输出与接地接线柱）和高频部分（频率调节、等幅调幅转换、幅度调节与输出接线柱）。这种分组方式与内部电路的划分一致，熟悉分组后即可从面板推知信号流向：低频振荡器的输出一路经幅度电位器送到低频输出端，另一路经转换开关送往高频部分的调幅管；高频振荡器的输出经幅度电位器送到高频输出端。理解这一点，操作时就不会把两路信号的调节旋钮混用。

下面按上述分组分别说明各控制件的作用。

\begin{enumerate}
\item \textbf{电源开关(S\(_2\))}：仪器电源开关，当此开关扳向“开”时，仪器接通电源，其左边发光二极管即发出红光，经预热后，仪器即可正常工作。
\item \textbf{0.2A(F\(_1\))}：仪器电源保险丝盒。
\item \textbf{\(\perp\)(X\(_5\))}：仪器接地端。
\end{enumerate}

低频部分

\begin{enumerate}
\item \textbf{频率kHz(S\(_1\))}：低频信号频率转换开关。扳动此开关可改变低频信号频率，分 0.5、1、1.5、2、2.5kHz 五档。
\item \textbf{幅度(RP\(_3\))}：低频信号输出幅度调节旋钮。顺时针旋转时，输出信号幅度加大。反时针旋转时，输出信号幅度减小，最后可调到零。仪器关机时应将其反时针旋转到底，使用时渐渐顺时针转动，使输出幅度达到要求即可。
\item \textbf{输出(X\(_1\))}：低频信号输出接线柱。使用时将连接电缆红色插头插入此接线柱。
\item \textbf{\(\perp\)(X\(_2\))}：低频信号输出接地接线柱。使用时将连接电缆黑色插头插入此接线柱。
\end{enumerate}

高频部分

\begin{enumerate}
\item \textbf{频率kHz(C\(_{21}\))}：仪器高频信号频率调节旋钮。转动此旋钮时，高频信号频率可由450kHz连续调节到1700kHz，由红色指针指出，其中465kHz为中频信号。
\item \textbf{等幅、调幅(S\(_3\))}：高频信号等幅波、调幅波转换开关。当拨向“等幅”，输出高频信号等幅波；当拨向“调幅”时，输出高频信号受低频信号调幅，调幅度由低频信号幅度调节旋钮控制，调幅频率由低频信号频率开关决定。
\item \textbf{幅度(RP\(_4\))}：高频信号输出幅度调节旋钮。顺时针旋转时，输出信号幅度加大。反时针旋转时，输出信号幅度减小，最后可调到零。仪器关机时应将其反时针旋转到底，使用时渐渐顺时针转动，使输出幅度达到要求即可。
\item \textbf{输出(X\(_3\))}：高频信号输出接线柱。使用时将连接电缆红色插头插入此接线柱。
\end{enumerate}

两个幅度调节旋钮都要求关机时反时针旋到底、使用时再渐渐顺时针旋出，这一规定有两重用意：一是避免开机瞬间大幅度信号突然加到被测电路上，二是使操作者养成"从零开始加量"的习惯，便于观察被测电路在小信号下的正常响应，防止过激励造成的饱和与误判。接地端的作用是把仪器外壳与被测设备的地连成一点，减小分布电容引入的干扰；若两台设备各自接地而地电位不同，反而可能引入交流声，因此实验中通常只在一处接地。

需要注意面板标注与内部记号的对应关系。同一功能的旋钮在不同的产品批次中，其元件记号可能不同，故后文电路原理与故障修理中出现的记号应以随机附带的电路图为准；本篇正文与元件表中的记号取自原说明书，阅读时应对照实物核对，不可仅凭记号臆断接线。

\section{电路原理}

J2465-2 型学生信号发生器由低频振荡器、高频振荡器及电源部分组成。其电路原理图见图 3。现将各部分作用原理分别叙述如下：

就总体框图而言，本仪器的信号流程可概括为：电源变压器降压、整流、滤波、稳压后向全机供电；低频振荡器产生五个点频的音频正弦信号，经射极跟随器与幅度电位器输出，同时可经转换开关送至调幅管；高频振荡器产生中波频段的正弦信号，其供电电压受调幅管控制，因而在"调幅"位置时输出被低频信号调制，最后同样经射极跟随器与幅度电位器输出。可以看出，本机把前言中所述"频率产生—电平控制—调制"三个环节以最简的形式实现：频率产生用 RC 与 LC 两种振荡器，电平控制用电位器分压，调制则采用改变振荡管供电电压的方式。
	
\subsection{低频振荡器}

低频振荡器采用 RC 桥式振荡电路，这种电路产生的波形失真小，频率稳定度高，特别适合产生低频信号。振荡器由 V\(_1\)、V\(_2\) 两只晶体管组成，V\(_3\) 为射极跟随器。为了说明其工作原理，振荡器电路简化如图1所示，可以看出振荡器的反馈网络是一个 RC 电桥，电桥的四个桥臂中，AB、BC 构成 RC 串并联选频网络，AD、DE 由负反馈电阻 R\(_3\)、R\(_6\) 构成。
	
	\begin{figure}[H]
		\centering
		% \includegraphics[width=0.8\textwidth]{rc_oscillator.png}
		\caption{低频振荡器电路原理图}
		\label{fig:low_freq_osc}
	\end{figure}
	
	在串并联选频网络中，\(R_1=R_2\)，\(C_1=C_2=C_3\)，当频率 \(f=1/(2\pi R_1C_1)\) 的信号通过网络时，网络输出电压的幅度为最大，是输入电压的 1/3，同时输出信号与输入信号相位相同，这个频率，就是电路的谐振频率。当频率大于或小于谐振频率时，网络输出电压的幅度都很小，这就达到了选频的目的，其选频特性如图 2 所示。
	
	\begin{figure}[H]
		\centering
		% \includegraphics[width=0.6\textwidth]{figure2.png}
		\caption{RC振荡器选频特性}
		\label{fig:rc_response}
	\end{figure}
	
	选频网络的输出信号，经两级放大后，输送到输入端，只要放大器放大倍数等于 3，电路即能自激振荡，放大器的放大量大一些，对振荡起振有利，但振荡波形失真。为了减小失真，应使电路工作于临界状态，即放大倍数正好为 3。这样就要求放大器工作很稳定，否则当工作点改变使放大倍数小于 3 时，电路就会停振。为此在放大器中引入了深的负反馈，以稳定放大器的工作状态。电桥 AD、DE 两臂的负反馈电阻就是起这个作用。调节 R\(_5\) 值，可使放大器放大倍数略大于 3，电路即能稳定地自激振荡。当 S\(_1\) 开关改变 R\(_1\)、C\(_1\)、R\(_2\)、C\(_2\) 的值时，振荡频率可在 500Hz、1kHz、1.5kHz、2kHz、2.5kHz 五个频率中变化。低频振荡器产生的信号，经 V\(_2\) 组成的射极跟随器输出，以增加振荡器带负载能力。RP\(_3\) 为输出信号幅度调节电位器，输出的低频信号还可以通过 S\(_4\) 开关加到高频信号调幅管，对高频信号进行调幅。
	
	从振荡理论看，任何自激振荡器都必须同时满足两个条件：相位条件要求环路总相移为 $2n\pi$，使反馈到输入端的信号与原信号同相；幅值条件要求环路增益的模不小于 1，即 $|\dot{A}\dot{F}| \ge 1$，其中 $\dot{A}$ 为放大器增益，$\dot{F}$ 为反馈系数。RC 串并联（文氏桥）网络在谐振频率上恰好满足相移为零，且此时传输系数为 $1/3$，即 $|\dot{F}| = 1/3$，故要求放大器的电压增益 $|\dot{A}| \ge 3$。谐振频率由网络元件决定
	\[
	f_0 = \frac{1}{2\pi R C},
	\]
	当 $R_1 = R_2 = R$、$C_1 = C_2 = C$ 时成立。这正是本机用波段开关同时切换 $R$、$C$ 数值以改变点频的依据：五个点频对应五组不同的 $RC$ 值。
	
	幅值条件说明了为什么必须引入深度负反馈。若增益远大于 3，振荡幅度将不断增长直至放大管进入非线性区，波形顶部被削平，失真度上升；若增益略小于 3，则振荡衰减而停振。因此实用电路总是把增益调整在略大于 3 的临界状态，并借助负反馈网络稳定增益：当幅度增大时负反馈加深、等效增益下降，从而把幅度稳定在某一数值上，这种自动稳幅作用是 RC 桥式振荡器失真小的根本原因。原电路中由 $R_3$、$R_6$ 构成的桥臂以及可调元件 $R_5$，正是用于设定这一临界增益。
	
	振荡器之后接射极跟随器，其作用是阻抗变换与隔离。射极跟随器的电压增益接近 1 而输入阻抗高、输出阻抗低，接在振荡器与负载之间，既避免外接负载的变化反过来影响选频网络（那将使频率和幅度都随负载漂移），又提高了带负载能力。理解这一环节的作用，也就理解了为什么输出幅度指标必须注明负载条件。
	
	\subsection{高频振荡器}
	高频振荡器由场效应管 V\(_6\) 组成电容三点式振荡电路。在三种 LC 振荡电路中，这种电路振荡频率最稳定，高频谐波最小，振荡波形最好。由于场效应管输入阻抗很高，对谐振电路影响很小，因而更进一步提高了振荡器的工作稳定性。振荡电路由电感 L\(_1\) 及双连可变电容 C\(_{11}\) 组成。双连可变电容两组串联，既是 LC 回路电容，又构成分压反馈电路。电路的振荡频率按下式计算：
	\[
	f = \frac{1}{2\pi \sqrt{\frac{1}{2} L_1 C_{11}}}
	\]
	
	当转动 C\(_{11}\) 双联可变电容时，高频信号的频率即可由 450kHz 到 1700kHz 连续变化。高频信号通过 V\(_6\) 组成的射极跟随器输出，RP\(_1\) 是输出幅度调节电位器。V\(_7\) 组成高频振荡器的调幅器，其射极输出电压作为振荡管 V\(_6\) 栅极电源。当调幅等幅开关 S\(_3\) 置“等幅”时，V\(_7\) 射极输出电压固定在 9 伏不变，振荡管产生高频等幅振荡。当 S\(_3\) 开关置“调幅”时，低频信号加到调幅管 V\(_7\) 基极，其发射极输出电压 9 伏上下按低频信号幅度被低频信号所调制。
	
	电容三点式（考毕兹）电路的特点是反馈电压取自回路电容的分压。设两只串联电容分别为 $C_a$、$C_b$，则回路的等效电容为
	\[
	C = \frac{C_a C_b}{C_a + C_b},
	\]
	振荡频率为 $f = 1/(2\pi\sqrt{L C})$；当两只电容相等（如本机采用双连可变电容两组串联）时 $C = C_a/2$，即得前式的形式。反馈系数近似为两电容之比
	\[
	F \approx \frac{C_a}{C_b},
	\]
	它决定了起振的难易与波形质量：反馈过弱不易起振，过强则波形失真、谐波增多。由于反馈取自电容分压，而电容对高次谐波呈现更小的阻抗，谐波成分被更多地旁路，故电容三点式电路输出波形的谐波含量较小，这正是它优于电感三点式电路之处。此外，采用场效应管作振荡管的意义在于其栅极输入阻抗极高，接入谐振回路时几乎不引入损耗，回路的品质因数得以保持，因而频率稳定度与波形都更好。
	
	改变振荡管的供电电压即可实现调幅，这是最简单的调幅方法之一：振荡幅度大体随供电电压变化，于是低频信号叠加在直流供电上时，高频输出的包络便按低频信号的规律起伏。所得已调波可写成
	\[
	u(t) = U_c\left[1 + m\cos(2\pi f_m t)\right]\cos(2\pi f_c t),
	\]
	式中 $U_c$ 为载波幅度，$f_c$ 为载波频率，$f_m$ 为调制频率，$m$ 为调幅度。把上式展开可知，已调波的频谱由载波 $f_c$ 与上下两个边频 $f_c + f_m$、$f_c - f_m$ 组成，占用带宽为
	\[
	B = 2 f_m .
	\]
	因此用 1~kHz 的低频信号调幅时，输出信号占用约 2~kHz 带宽，检修接收机中频通道时可据此判断通带是否足够。使用中还须注意：调幅度由低频幅度旋钮控制，旋得过大将进入过调制而使包络底部削平，输出出现明显失真；同时由于调幅是通过改变振荡管供电实现的，调制过深也会牵动振荡频率，产生附带的寄生调频。原设计采用漏极（供电端）调幅并配以射极跟随器输出，就是为了在获得足够调幅度的同时把载频偏移限制在可接受的范围内。
	
	\subsection{电源部分}
	电源部分由电源变压器、整流滤波电路及稳压电路组成。220 伏交流电经 T\(_1\) 变压器降压，在 3、4、5 端输出 18 伏交流电压，经 V\(_8\)、V\(_9\) 全波整流、C\(_{16}\)、C\(_{17}\)、R\(_{20}\) 组成的 \(\pi\) 型滤波网络滤波，再由稳压管 V\(_{12}\) 稳压后，供给电路各部分。V\(_{10}\) 为发光二极管作为仪器指示灯。当仪器接通电源，电源部分正常工作时，即发出红光。
	
	电源电路的四个环节各有明确的定量关系。变压器把电网电压降到所需的交流值 $U_2$；全波整流把交流变为脉动直流，在电阻负载下其平均值为
	\[
	U_{av} = \frac{2\sqrt{2}}{\pi} U_2 \approx 0.9 U_2 ,
	\]
	而在接有较大滤波电容时，输出直流电压将接近变压器副边电压的峰值 $\sqrt{2}\,U_2$，这就是 18~V 交流电经整流滤波后可得到十六至二十余伏直流的原因。滤波环节利用电容对交流的低阻抗和电阻对交流的分压作用抑制纹波，本机采用的 $\pi$ 型（电容—电阻—电容）滤波网络对纹波的衰减倍数近似为
	\[
	K \approx \omega C_2 R ,
	\]
	式中 $\omega$ 为纹波角频率（全波整流时为电网频率的两倍），可见增大电容或电阻都能改善滤波效果，但增大串联电阻会加大直流压降，故两者需折中。
	
	稳压环节采用稳压二极管并联稳压，其工作要点是稳压管必须始终处于反向击穿状态。设整流滤波后的电压为 $U_i$、稳压值为 $U_z$、负载电流为 $I_L$，则限流电阻应满足
	\[
	R = \frac{U_i - U_z}{I_z + I_L},
	\]
	且要求在电网电压最低、负载电流最大时稳压管电流仍不低于其最小稳定电流，在电网电压最高、负载电流最小时不超过其最大允许电流。稳压的意义对本机尤其重要：振荡器的频率与幅度都随供电电压变化，供电稳定是频率稳定度和幅度稳定的前提。发光二极管作指示灯时必须串联限流电阻，其阻值按
	\[
	R = \frac{U - U_F}{I_F}
	\]
	估算，式中 $U_F$、$I_F$ 分别为发光二极管的正向电压与工作电流。指示灯直接取自稳压后的电源，故它不亮往往意味着电源部分有故障，这一现象在故障判断中很有价值。
	
	\begin{figure}[H]
		\centering
		% \includegraphics[width=0.9\textwidth]{figure3.png}
		\caption{J2465-2型学生信号发生器电原理图}
		\label{fig:schematic}
	\end{figure}
	
	\begin{figure}[H]
		\centering
		% \includegraphics[width=0.9\textwidth]{figure4.png}
		\caption{J2465-2型学生信号发生器内部结构示意图}
		\label{fig:structure}
	\end{figure}
	
	\begin{figure}[H]
		\centering
		% \includegraphics[width=0.9\textwidth]{figure5.png}
		\caption{J2465-2型学生信号源印制电路板}
		\label{fig:pcb}
	\end{figure}
	
	\section{使用方法}
	\begin{enumerate}
		\item 仪器使用前，应仔细阅读本说明书，以帮助你了解仪器的工作原理，使用方法，防止因使用不当而损坏仪器。
		\item 将仪器插上 220 伏 \(\pm 10\%\) 电源，打开电源开关，其左边指示灯即发光，经预热后仪器即可使用。
		\item 需要低频信号输出时，将电缆连接线两只插头接到低频信号输出接线柱 X\(_4\) 及地线接线柱 X\(_2\)，电缆线的另一端接到测试线路。将频率开关 S\(_1\) 拨到需要的频率数，调幅、等幅开关 S\(_2\) 置“等幅”，渐渐顺时针转动幅度旋钮 (RP\(_3\))，即有低频信号输出。
		\item 需要高频信号输出时，将电缆连接两只插头接到高频信号输出接线柱 X\(_3\) 及地线接线柱 X\(_2\)。转动高频频率调节旋钮 RP\(_2\)，使指针指到需要的频率值；渐渐顺时针转动幅度旋钮 RP\(_1\)，即有高频信号输出。如果需要等幅波信号，则将等幅、调幅开关 S\(_3\) 置“等幅”；如果需要调幅波信号，则将等幅、调幅开关 S\(_3\) 置“调幅”，再渐渐顺时针转动低频信号幅度旋钮 RP\(_3\)，高频信号即被低频信号调幅。顺时针转动角度越大，高频信号的调幅度越深。
		\item 仪器高频信号和低频信号可以同时输出，但如果只用一种信号时，最好将另一信号连接电缆拔去，并将幅度旋钮反时针旋到底。
		\item 测试时高频信号和低频信号输出有足够幅度即可，没有必要时不要将幅度旋钮顺时针旋到底。
		\item 高频信号和低频信号输出时，测试线路连接点输入电阻不应大于 300 欧姆，以免影响仪器正常工作，并特别防止将输出电缆两头短路，以造成仪器损坏。
		\item 仪器使用完毕后，将高频信号幅度旋钮及低频信号幅度旋钮均反时针旋到底，关掉仪器电源开关，并拔出电源插头，整理好电源线，将仪器放到阴凉干燥通风的仪器架上。
		\item 仪器存放满三个月，必须开机一次，开机时间不短于 2 小时。
	\end{enumerate}
	
	上列操作规程中有几条值得进一步说明其道理。要求被测电路输入电阻不低于 $300\,\Omega$，是因为输出级的带负载能力有限，负载过轻（阻值过小）会使输出幅度按前述分压关系明显下降，并可能因输出管电流过大而失真；输出端短路则相当于负载为零，射极跟随器将承受远超正常值的电流，极易损坏。要求"没有必要时不要把幅度旋到底"，一方面是避免输出级在大幅度下工作时的失真增大，另一方面是保护被测电路——检修接收机时若注入信号过强，各级都将饱和，无法判断哪一级增益不足。
	
	用本仪器检修收音机时，通常采用信号注入法，即由后级向前级逐点注入信号并监听输出。具体做法是：先在检波器之后（音频输入端）注入低频信号，检查音频放大与扬声器通道；再在中频放大器输入端注入 465~kHz 调幅信号，检查中放与检波；最后在天线输入端注入相应频率的调幅信号，检查输入回路与变频级。每前移一级，若输出声音消失或明显减弱，则故障多在刚跨过的那一级。注入时应通过隔直电容或按仪器规定的连接电缆接入，避免把被测电路的直流电位引到仪器输出端；同时把仪器接地端与被测机底板可靠连接，否则分布电容引入的交流声会干扰判断。
	
	关于存放与定期通电，其原因在于元件的物理特性：铝电解电容长期不通电时氧化膜会逐渐变薄，漏电流增大，通电运行一段时间可使氧化膜自行修复；机内受潮会使绝缘电阻下降、可变电容器与开关触点氧化，通电产生的少量热量有利于驱除潮气。因此"存放满三个月开机一次、每次不短于两小时"并非可有可无的建议，而是保证仪器长期可用的必要维护。此外，可变电容器与波段开关是最易磨损的机械部件，操作时应缓慢均匀转动，切勿快速反复扳动或用力过猛。
	
	\section{仪器故障修理}
	
	检修本仪器宜遵循由外到内、由粗到细的一般程序。第一步做直观检查：观察保险丝、电源线、插头与接线柱是否完好，机内有无元件烧焦、变色、脱焊和引线相碰。第二步测量电源：确认电网电压在允许范围内，再测整流滤波后与稳压后的直流电压是否正常，电源不正常时机内所有故障现象都不足为凭。第三步分段判断：本机由低频、高频两个独立通道与公用电源组成，据故障现象即可先定通道——两路都无输出多为电源问题，仅一路无输出则故障在该通道。第四步在可疑通道内用电压法测各级静态工作点、用示波器逐点观察波形，找出信号消失或波形异常的位置。第五步对可疑元件采用替代法或断开测量法确认，特别是晶体管、场效应管与电解电容，其失效往往无法从外观判断。
	
	下表按故障现象归纳了可能的部位与相应的检查手段，可作为检修时的索引；各项的具体分析见以下各小节。
	
	\begin{table}[H]
		\centering
		\caption{常见故障现象与检查手段索引}
		\begin{tabular}{|p{3.2cm}|p{4.4cm}|p{4.8cm}|}
			\hline
			\textbf{故障现象} & \textbf{可能部位} & \textbf{检查手段} \\ \hline
			全机无输出、指示灯不亮 & 保险丝、变压器、整流滤波、稳压 & 测电网电压、整流与稳压输出电压 \\ \hline
			低频某档无输出 & 该档选频网络元件变值或虚焊 & 测该档电阻电容、观察起振情况 \\ \hline
			低频各档均无输出 & 振荡器工作点偏移、管子失效 & 测供电电压与静态工作点，调可调元件 \\ \hline
			低频输出失真 & 增益过大、跟随器损坏 & 示波器观察波形，测跟随器输出 \\ \hline
			高频全频段无输出 & 回路电感断路、双连电容接地不良 & 通断测量、测振荡管电极电压、示波器观察 \\ \hline
			高频局部频段无输出 & 双连电容器动静片碰片 & 缓慢转动观察、目视检查片间 \\ \hline
			调幅波失真 & 调幅管或振荡管工作点偏移 & 测调幅管射极电压、示波器观察包络 \\ \hline
		\end{tabular}
	\end{table}
	\subsection{仪器无低频信号输出}
	如果某一档无输出，则表示 RC 桥式振荡电路在该档停振，其原因是选频网络 R\(_1\)、R\(_2\)、C\(_1\)、C\(_2\)、R\(_3\)、R\(_4\)、R\(_5\)、R\(_{20}\) 中，该档对应的电阻量变小，可以在与它相串联的电阻 R\(_{11}\)、R\(_{12}\)、R\(_{13}\)、R\(_{14}\)、R\(_{15}\)、R\(_{16}\)、R\(_{17}\)、R\(_{18}\)、R\(_{19}\)、R\(_{20}\) 中所对应的电容上并联一只大电阻，如果电阻值合适，即能起振，并满足幅度要求。如果所有档都无输出，其主要原因可能是 RC 桥式振荡电路工作点偏移引起停振。修理时首先检查电网电压是否在 220V \(\pm 10\%\) 范围内，其次检查在振荡器直流供电电压是否在 16--18V 范围内，若都正常，可调节 RP\(_2\)，使其起振，波形输出最大且不失真。若调节 RP\(_2\) 也不能起振，则要检查 C\(_{21}\)～C\(_{23}\) 是否虚焊，V\(_1\)、V\(_2\) 管是否损坏，电阻是否失效，射极跟随器 V\(_3\) 是否正常工作，故障排除后再调 RP\(_2\)，使其起振。
	
	判断这类故障时应把握"停振"的实质：环路增益降到 1 以下，或相位条件被破坏。前者对应放大器增益不足（工作点偏移、管子放大倍数下降、负反馈过深），后者对应选频网络元件变值或断路。因此检修顺序宜为：先确认供电电压正常，再用示波器观察放大级有无自激的迹象，最后核对该档选频元件。调整可调元件使其起振后，应用示波器复核波形不失真、并检查其余各档是否仍能起振，避免为使某一档起振而把增益调得过大，导致其他档失真。
	
	\subsection{低频信号输出失真}
	如果某一档失真，其原因是选频网络中 R\(_1\)、R\(_2\)、R\(_3\)、R\(_{11}\)、R\(_{12}\)、R\(_{13}\)、R\(_{14}\) 中，该档所对应的电阻量变大，可以在与它相串联的电阻 R\(_{11}\)、R\(_{12}\)、R\(_{13}\)、R\(_{14}\)、R\(_{15}\)、R\(_{16}\)、R\(_{17}\)、R\(_{18}\)、R\(_{19}\)、R\(_{20}\) 中所对应的电阻上并联一只大电阻，如电阻合适，即可消除失真。如果各档输出低频都失真，其原因为 RC 桥式振荡器工作点偏移引起的，可适当调节 RP\(_2\)，使输出波形消除失真，并保证有足够大的输出。如果输出低频是半波信号，则表示射极跟随器 V\(_3\) 损坏或偏置电阻 R\(_{20}\) 脱开。
	
	失真的形态对判断部位很有帮助，检修时应先用示波器看清波形而不是急于更换元件。波形顶部与底部同时被对称削平，多为振荡幅度过大、放大器进入非线性区，属增益偏高问题；仅单侧削平或呈半波，则说明某一级工作点严重偏移或跟随器失效，属直流偏置问题；波形上叠加有明显的低频起伏或毛刺，应检查滤波电容是否失效、有无外界干扰窜入。测量失真度时还须注意，若被测点后面接有幅度电位器，电位器接触不良同样会造成波形畸变，可旋动旋钮观察波形是否随之抖动来加以区分。
	
	\subsection{仪器无高频信号输出}
	如果在整个频段都无高频信号输出，最常见原因是高频振荡器中电感线圈 L\(_1\) 断路或焊点脱开，使电容之点或振荡器中失去工作电压。也可能是双联电容器外壳接地不良。如上述两种原因均被排除，则可用万用表检查 V\(_6\) 管射极电压是否为正常 8 伏左右，如不正常，则表示 V\(_6\) 管损坏或 R\(_{31}\)、R\(_{32}\) 电阻脱焊。如果电压正常，则可用示波器观察 V\(_6\) 栅极输入处有没有高频信号，如没有，则可能是场效应管 V\(_8\) 失效或电阻 R\(_{53}\)、R\(_{54}\)、电容 C\(_6\)、C\(_{22}\)、C\(_{24}\) 损坏。如有高频信号，则故障在射极跟随器 V\(_{10}\) 及输出电路。
	
	高频电路的检修比低频电路更依赖正确的测量方法。用普通万用表在高频回路上直接测量交流电压是不可靠的，表笔引入的分布电容足以使振荡停止或频率大幅偏移，故观察高频信号应使用带高频探头的示波器或高频电压表，测量点尽量选在射极跟随器输出等低阻抗处。同时，谐振回路的检修要格外注意机械因素：电感线圈的焊点、双连可变电容的转轴接地弹片与外壳的接触，都可能因氧化或松动而时断时通，表现为敲击机壳时输出忽有忽无，这类故障用静态测量往往查不出，需在通电状态下轻轻拨动可疑部位观察输出变化。
	
	\subsection{高频信号输出失真}
	如果在高频信号频段中某一部分无信号输出，其原因为双联电容器动片和静片局部碰片，可以将碰片部位拨开。
	
	碰片的判断方法是：断电后缓慢转动调谐旋钮，用万用表电阻挡监视双连电容动片与静片之间，正常时应始终为断路，某一角度出现短路即为碰片位置。拨正时应用无磁性的薄片工具轻轻分开，切勿用力过大以致片子变形，否则将造成容量变化、频率刻度整体偏移。若某一频段输出幅度明显低于其他频段但并未消失，则不一定是碰片，还可能是回路在该频段的品质因数下降（如线圈受潮、磁芯松动），可结合幅度随频率变化的规律加以区分。
	
	\subsection{高频调幅波失真}
	其主要原因是场效应振荡管 V\(_6\) 工作点偏移，可调换 R\(_{31}\)、R\(_{32}\) 电阻值使其正常或更换场效应管。
	
	检修前应先排除"操作性失真"，即把低频幅度旋钮回到较小位置重新观察：若包络底部削平随旋钮加大而出现，则属过调制，并非故障。确属故障时，可用示波器同时观察低频信号与已调波包络：低频信号本身失真则应先修低频通道；低频正常而包络上下不对称、或包络最小处已到零，说明调幅管的静态射极电压偏离设计值，使调制工作点不在线性段，此时按原文调整偏置电阻或更换管子即可。此外，若调幅时载频明显偏移（可用频率计监视），多为调幅深度过大或振荡管工作点不当所致。
	
	\subsection{指示灯不亮}
	如果高频低频信号均有输出，则表示只是指示灯损坏，或 R\(_{22}\) 电阻脱开。如高低频信号均无输出，则表示仪器电源部分有故障，可检查电源保险丝是否完好，稳压管 V\(_{12}\) 是否击穿，印制板电源部分线路是否和变压器双连电容器相碰。
	
	指示灯的价值在于它是电源的"活标志"，因此该现象须与其他现象合起来判断：灯不亮而两路信号正常，问题只在指示支路（发光二极管本身或其限流电阻）；灯不亮且两路信号皆无，则应按电源电路顺序检查保险丝、变压器副边电压、整流二极管、滤波电容与稳压管。检查中若发现保险丝反复熔断，切不可换用更大规格的保险丝了事，而应查明短路点，常见者为滤波电容击穿、稳压管短路或印制板受潮爬电。更换发光二极管时须注意极性与限流电阻的匹配，不可省去限流电阻直接接入电源。
	
	最后需要提醒的是安全事项。机内变压器初级及电源引线部分带 220~V 交流电，通电检修时必须避免触及，测量前应确认表笔绝缘完好；断电后滤波电容上仍可能存有电荷，应先泄放再动手。焊接印制板时宜使用功率适中、带接地的电烙铁，防止漏电与静电损坏场效应管；拆装场效应管时应做好人体与工具的静电防护。
	
	\appendix
	\section{J2465-2型学生信号发生器元件表（一）}
	
	元件表按电路记号列出各元件的型号规格，检修与配件采购时应以此为依据。表中的型号采用旧国家标准的命名方式，其构成一般是"类别字母+材料或结构字母+序号+主要参数+允许偏差"。半导体器件的型号中，第一位数字表示电极数目（2 为二极管、3 为三极管），第二位字母表示材料与极性，第三位字母表示器件类型，例如 Z 表示整流管、W 表示稳压管、G 表示高频小功率管、J 表示结型场效应管；开关、接线柱等结构件则以厂标代号表示其形式与规格。代换元件时应遵循"参数不低于原件、类型相同"的原则：整流与稳压器件须核对反向耐压、正向电流或稳压值与耗散功率，晶体管须核对特征频率、集电极最大电流与耗散功率，场效应管还须注意跨导与结构类型是否一致。
	
	\begin{table}[H]
		\centering
		\caption{晶体管与开关等元器件表}
		\begin{tabularx}{\textwidth}{|c|X|}
			\hline
			\textbf{电路记号} & \textbf{型号规格} \\ \hline
			RP\(_2\) & WH5-1-1W-1-5k-X-202S-1 \\ \hline
			S\(_1\) & KCX-5W2D \\ \hline
			S\(_2\) & KNX-2W1D \\ \hline
			S\(_3\) & KB05-2W2D \\ \hline
			V\(_1\) & 3DG100A \\ \hline
			V\(_2\) & 3DG100A \\ \hline
			V\(_3\) & 3DG100A \\ \hline
			V\(_4\) & FG323001 \\ \hline
			V\(_5\) & 2CZ82C \\ \hline
			V\(_6\) & 2CZ82C \\ \hline
			V\(_7\) & 2CW113 \\ \hline
			V\(_8\) & 3DG100A \\ \hline
			V\(_9\) & 3DJ7G \\ \hline
			V\(_{10}\) & 3DG100A \\ \hline
			X\(_1\) & 接线柱910A \\ \hline
			X\(_2\) & 接线柱910A \\ \hline
			X\(_3\) & 电源插头 \\ \hline
			X\(_4\) & 接线柱910A \\ \hline
			L\(_1\) & 电感线圈 \\ \hline
			T\(_1\) & 变压器 \\ \hline
			F\(_1\) & 保险丝 BGXP0.2A \\ \hline
		\end{tabularx}
	\end{table}
	
	\section{J2465-2型学生信号发生器元件表（二）}
	由于元件较多，表二分为电阻和电容两部分列出。
	
	\subsection{电阻部分}
	
	电阻器型号中，首位字母 R 表示电阻器，第二位字母表示电阻体材料，例如 T 为碳膜、J 为金属膜；其后依次为结构或系列代号、额定功率（单位为瓦）、标称阻值和允许偏差或精度等级。表中标注 $\pm 1\%$ 的电阻多用于选频网络与分压网络等对精度要求较高之处，检修时不可用普通等级电阻随意代换，否则将直接影响振荡频率或输出幅度的准确性。其余通用位置的电阻，代换时阻值应尽量取原值，额定功率不得低于原件；电位器则以 W 打头，第二位字母表示电阻体材料（如 H 为合成碳膜），代换时须同时核对阻值、功率、轴长与阻值变化规律（线性或对数）。
	
	\begin{table}[H]
		\centering
		\small
		\setlength{\tabcolsep}{4pt}
		\caption{电阻元件表}
		\begin{tabular}{|c|c||c|c||c|c|}
			\hline
			\textbf{记号} & \textbf{规格} & \textbf{记号} & \textbf{规格} & \textbf{记号} & \textbf{规格} \\ \hline
			R\(_1\) & RTX-0.125-150-Ⅰ & R\(_{16}\) & RTX-0.125-68k-Ⅰ & R\(_{31}\) & RTX-0.125-3k-Ⅰ \\ \hline
			R\(_2\) & RTX-0.125-1.5k-Ⅰ & R\(_{17}\) & RTX-0.125-200k-Ⅰ & R\(_{32}\) & RTX-0.125-3k-Ⅰ \\ \hline
			R\(_3\) & RTX-0.125-1M-Ⅰ & R\(_{18}\) & RTX-0.125-10k-Ⅰ & R\(_{33}\) & RTX-0.125-1k-Ⅰ \\ \hline
			R\(_4\) & RTX-0.125-10k-Ⅰ & R\(_{19}\) & RTX-0.125-1k-Ⅰ & R\(_{34}\) & RTX-0.125-200k-Ⅰ \\ \hline
			R\(_5\) & RTX-0.125-10k-Ⅰ & R\(_{20}\) & RTX-0.125-200k-Ⅰ & R\(_{35}\) & RTX-0.125-3k-Ⅰ \\ \hline
			R\(_6\) & RTX-0.125-5.1k-Ⅰ & R\(_{21}\) & RTX-0.125-10k-Ⅰ & R\(_{36}\) & RTX-0.125-3k-Ⅰ \\ \hline
			R\(_7\) & RTX-0.125-100-Ⅰ & R\(_{22}\) & RJ-0.125-14.5k\(\pm\)1\% & R\(_{37}\) & RTX-0.125-3k-Ⅰ \\ \hline
			R\(_8\) & RTX-0.125-470-Ⅰ & R\(_{23}\) & RJ-0.125-8.2k\(\pm\)1\% & R\(_{38}\) & RTX-0.125-3k-Ⅰ \\ \hline
			R\(_9\) & RTX-0.125-20k-Ⅰ & R\(_{24}\) & RJ-0.125-5.6k\(\pm\)1\% & R\(_{39}\) & RTX-0.125-3k-Ⅰ \\ \hline
			R\(_{10}\) & RTX-0.125-200k-Ⅰ & R\(_{25}\) & RTX-0.125-1k-Ⅰ & R\(_{40}\) & RTX-0.125-3k-Ⅰ \\ \hline
			R\(_{11}\) & RTX-0.125-1.5k-Ⅰ & R\(_{26}\) & RJ-0.125-10k\(\pm\)1\% & R\(_{41}\) & RTX-0.125-3k-Ⅰ \\ \hline
			R\(_{12}\) & RJ-0.5-47-Ⅰ & R\(_{27}\) & RJ-0.125-10k\(\pm\)1\% & R\(_{42}\) & RTX-0.125-3k-Ⅰ \\ \hline
			R\(_{13}\) & RTX-0.125-5.1k-Ⅰ & R\(_{28}\) & RTX-0.125-100-Ⅰ & R\(_{43}\) & RTX-0.125-3k-Ⅰ \\ \hline
			R\(_{14}\) & RTX-0.125-1k-Ⅰ & R\(_{29}\) & RTX-0.125-200k-Ⅰ & R\(_{44}\) & RTX-0.125-3k-Ⅰ \\ \hline
			R\(_{15}\) & RTX-0.125-10k-Ⅰ & R\(_{30}\) & RTX-0.125-15k-Ⅰ & & \\ \hline
		\end{tabular}
	\end{table}
	
	\subsection{电容与电位器部分}
	
	电容器型号中，首位字母 C 表示电容器，第二位字母表示介质材料，例如 C 为瓷介、L 为有机薄膜（涤纶）、D 为铝电解；其后为结构序号、额定电压和标称容量及允许偏差。不同介质的电容适用场合差别很大：瓷介电容体积小、高频损耗低，适用于高频回路与旁路；有机薄膜电容稳定性好、损耗小，本机选频网络中即采用此类；铝电解电容容量大但有极性、高频特性差、寿命受温度影响，适用于电源滤波与低频耦合、旁路。代换时应保证介质类型相当、耐压不低于原值、容量在允许偏差内，电解电容尤须注意极性不得接反。
	
	本机可变电容器与电位器属机械可调件，除电气参数外还须核对结构尺寸与安装方式；双连可变电容的两组容量应基本一致，否则将影响频率刻度的对应关系。表中同一记号在不同批次可能对应不同规格，配件时应以实物与随机电路图核对为准。
	
	\begin{table}[H]
		\centering
		\small
		\setlength{\tabcolsep}{4pt}
		\caption{电容与电位器元件表}
		\begin{tabular}{|c|c||c|c|}
			\hline
			\textbf{记号} & \textbf{规格} & \textbf{记号} & \textbf{规格} \\ \hline
			C\(_1\) & CL\(_{11}\)-63-22n\(\pm\)5\% & C\(_{15}\) & CCX-Y-b-X-100p\(\pm\)10\% \\ \hline
			C\(_2\) & CL\(_{11}\)-63-22n\(\pm\)5\% & C\(_{16}\) & CZJ1-160V-0.01\(\pm\)10\% \\ \hline
			C\(_3\) & CD\(_{11}\)-25V-10\(\mu\) & RP\(_1\) & WH06-2-1k \\ \hline
			C\(_4\) & CD\(_{11}\)-25V-10\(\mu\) & RP\(_2\) & WH06-2-1k \\ \hline
			C\(_5\) & CD\(_{11}\)-25V-10\(\mu\) & RP\(_3\) & WH5-1-1W-1-5k-X-202S-1 \\ \hline
			C\(_6\) & CD\(_{11}\)-25V-4.7\(\mu\) & & \\ \hline
			C\(_7\) & CD\(_{11}\)-25V-10\(\mu\) & & \\ \hline
			C\(_8\) & CD\(_{11}\)-25V-220\(\mu\) & & \\ \hline
			C\(_9\) & CD\(_{11}\)-25V-10\(\mu\) & & \\ \hline
			C\(_{10}\) & CDM-202B & & \\ \hline
			C\(_{11}\) & CCX-Y-b-1-1000p\(\pm\)10\% & & \\ \hline
			C\(_{12}\) & CCX-Y-b-II-30p\(\pm\)10\% & & \\ \hline
		\end{tabular}
	\end{table}

\backmatter

\end{document}